% vim: set tw=78 sts=2 sw=2 ts=8 aw et ai:

\subsection*{Results}

I run my new implementation on the same \href{https://github.com/ReGeLePuMa/testing}{dataset} as
with the original one and I have found some interesting results.

While the novel dynamic analysis technique demonstrated measurable improvements
in dependency detection compared to the traditional baseline, 
these gains were largely restricted to applications with locally accessible source code.
Furthermore, the integration of the AI model introduced significant operational limitations.
Notably, the model frequently generated malformed inputs that induced container crashes, despite
the attempts to mitigate this issues by modifying the system prompt, which consequently bypassed 
the primary analysis phase and triggered a fallback to the binary search brute-force. 
Moreover, the system exhibited an excessively high rate of token consumption, 
particularly during the evaluation of available source code.

\subsection*{Further Work}

Considering these results, overcoming the above-mentioned deficiencies is going to be 
the first priority in further stages, especially with regard to eliminating the issue of 
container crashes.

In this sense, areas of interest for further research will involve:
\begin{itemize}
    \item \textbf{Support for more models}: I want to add support for more providers, especially OpenAI and Anthropic
    \item \textbf{More dynamic analysis traces}: I want to add more probes to trace the application, including \textit{ftrace}, \textit{fanotify} etc.
    \item \textbf{Compose support}: I want to add a wrapper for docker compose, to minimize a whole stack at once, instead of doing it for each image separately.
\end{itemize}
